\documentclass[journal,compsoc]{IEEEtran}
\usepackage{cite}
\usepackage{amsmath,amssymb,amsfonts}
\usepackage{algorithmic}
\usepackage{graphicx}
\usepackage{textcomp}
\usepackage{xcolor}
\usepackage{booktabs}
\usepackage{url}
\usepackage{array}

\begin{document}

\title{A Privacy-Preserving and Heterogeneity-Aware Federated Learning Framework for Prostate Cancer Staging and Client Contribution Valuation}

\author{Anonymous Authors
\thanks{Affiliation anonymized for peer review.}
}

\maketitle

\begin{abstract}
Collaborative machine learning in medical networks is often hindered by strict data privacy regulations and statistical heterogeneity across clinical sites. While federated learning (FL) enables multi-institutional training without direct data aggregation, the multi-dimensional interactions between statistical heterogeneity, privacy constraints, local adaptation, and participant valuation remain poorly characterized. In this study, we systematically investigate the central research question: How do privacy constraints and statistical heterogeneity jointly affect federated clinical model utility, local adaptation, and participant contribution valuation under decentralized training? We evaluate this using a multi-dimensional framework, DP-FPS, which integrates FedProx optimization, sample-level differential privacy calibrated via Rényi Differential Privacy (RDP) under a full sampling rate $q=1.0$, personalized local adaptation (PFL), and game-theoretic permutation Shapley values. Using the TCGA-PRAD clinical cohort (572 sample records corresponding to 500 unique patients), our evaluations show that: (1) FedProx consistently constrains client weight drift by 4.2--4.3\% across all tested Dirichlet concentrations; (2) FedProx shows a positive utility difference relative to FedAvg across the evaluated privacy--heterogeneity configurations over repeated trials; (3) personalized local adaptation improves performance over local-only training in the evaluable client setting (+7.87\%), although it does not outperform the global FedProx model; (4) privacy noise perturbation altered the observed contribution rankings in the evaluated three-client setting; and (5) RDP-calibrated noise reduces membership inference attack (MIA) advantage, indicating substantially reduced vulnerability to the evaluated confidence-threshold MIA. Finally, we report that while a federated Cox proportional-hazards survival model was successfully developed, out-of-sample survival discrimination could not be reliably evaluated because the held-out test partition contained zero observed death events due to severe right-censoring (12 events overall in 347 matched samples).
\end{abstract}

\begin{IEEEkeywords}
Federated Learning, Differential Privacy, FedProx, Non-IID Learning, Prostate Cancer, Data Valuation, Shapley Value, Membership Inference.
\end{IEEEkeywords}

\section{Introduction}
\IEEEPARstart{P}{athologic} staging of prostate cancer is vital for clinical risk stratification and treatment planning. The training of robust staging models relies on large-scale datasets, yet centralizing clinical and genomic cohorts across hospitals is challenging due to strict regulatory compliance standards (e.g., HIPAA, GDPR). Federated Learning (FL) offers a collaborative paradigm by keeping raw data local and aggregating only model updates. 

However, cross-silo clinical FL faces three fundamental constraints: (1) \textit{Statistical Heterogeneity}: Differences in clinical protocols and demographics across hospitals yield non-IID client distributions, leading to optimization drift. (2) \textit{Privacy Vulnerability}: Exposing model updates to a coordinator introduces reconstruction risks, requiring differential privacy (DP) guarantees. (3) \textit{Consortium Stability}: Participating clients require fair incentives, necessitating stable contribution valuation.

Prior studies have typically evaluated these features as independent components. In contrast, this paper systematically analyzes the multi-dimensional trade-offs and joint interactions between these mechanisms. Our work is organized around a primary research question: **How do privacy constraints and statistical heterogeneity jointly affect federated clinical model utility, local adaptation, and participant contribution valuation?**

To address this question, we evaluate a multi-layered framework (DP-FPS) on the TCGA-PRAD clinical-protein cohort, establishing the following core contributions:
\begin{itemize}
    \item **Heterogeneity Bounding (H1)**: We demonstrate that FedProx consistently reduces client-update divergence by ~4.3\% across increasingly heterogeneous Dirichlet partitions.
    \item **Privacy-Utility Trade-offs (H2)**: Stronger differential privacy constraints (smaller composed $\epsilon$) degrade predictive utility, while moderate budgets preserve performance.
    \item **Privacy $\times$ Heterogeneity Interaction (H3)**: We show that FedProx consistently improves utility relative to FedAvg under joint constraints across repeated trials ($\mu_{\Delta} > 0$ with 95\% confidence intervals).
    \item **Personalized FL Adaptation (H4)**: Personalized local adaptation (PFL) improves performance over local-only training in the evaluable client setting, although it does not outperform the global FedProx model in this experiment.
    \item **Contribution Valuation under Privacy Noise (H5)**: We reveal that differential privacy noise altered the observed contribution rankings in the evaluated three-client setting, highlighting a critical tension between privacy and economic equity in collaborative networks.
    \item **Controlled Domain-Shift Robustness (H6)**: We characterize different model degradation profiles under controlled synthetic covariate and concept shifts.
    \item **Empirical Privacy Protection (H7)**: We execute a Membership Inference Attack (MIA) to show that RDP-calibrated noise reduces empirical vulnerability to the evaluated confidence-threshold attack.
\end{itemize}

\section{Related Work}
Collaborative clinical machine learning has historically relied on manual data transfer agreements or centralized registry models. Early federated medical studies applied FedAvg \cite{McMahan2017FedAvg} to clinical classification tasks but identified significant performance degradation when client populations exhibited clinical subpopulation skew. To mitigate this heterogeneity, Li et al. proposed FedProx \cite{Li2020FedProx}, adding a proximal regularization term to keep local updates close to the global consensus. 

Privacy preservation in medical FL has increasingly leveraged Differential Privacy (DP) \cite{Abadi2016DPSGD, Geyer2017PrivateFL}. Standard DP-SGD applies gradient clipping and Gaussian noise addition to client updates. Composing privacy losses over multiple rounds requires advanced accountants. Mironov et al. introduced Rényi Differential Privacy (RDP) \cite{Mironov2017RDP, Mironov2019SampledGaussian}, providing tight analytical composition bounds for the Gaussian mechanism, which is critical when training models over multiple local epochs.

Incentivizing hospital participation in federated networks remains an active area of research \cite{Kazlouski2025HospitalParticipation, Wang2020ContributionSV}. Leave-One-Out (LOO) methods evaluate a client's marginal value by comparing model performance with and without that client \cite{Ghorbani2019DataShapley}. However, LOO fails to capture multi-client synergies and exhibits high variance. Cooperative game theory provides the Shapley value \cite{Jia2019KNNShapley}, offering a fair payout distribution based on axiomatic principles, though approximating it under communication constraints is non-trivial.

\section{System Model / Problem Formulation}
We model a federated network of $K$ hospital clients and a central coordinator.
Each client $k \in \{1, \dots, K\}$ possesses a private local clinical-protein dataset $D_k = \{(x_{i,k}, y_{i,k})\}_{i=1}^{n_k}$, where $x_{i,k} \in \mathbb{R}^d$ represents a multimodal feature vector (clinical attributes and genomic components) and $y_{i,k} \in \{0, 1\}$ represents the pathologic staging target. Let $n = \sum_{k=1}^K n_k$ denote the aggregate size of the dataset.

The decentralized objective is to solve:
\begin{equation}
\min_w F(w) = \sum_{k=1}^K \frac{n_k}{n} F_k(w)
\end{equation}
where $F_k(w) = \frac{1}{n_k} \sum_{i=1}^{n_k} \ell(w; x_{i,k}, y_{i,k})$ is the local empirical risk of client $k$. 
Under FedAvg, clients perform local gradient updates for $E$ epochs and upload their updated parameters to the server. The coordinator averages the local updates:
\begin{equation}
w^{t+1} = \sum_{k=1}^K \frac{n_k}{n} w_k^{t+1}
\end{equation}
Under FedProx, each local objective is regularized to prevent client parameters from drifting away from the global consensus $w^t$:
\begin{equation}
h_k(w; w^t) = F_k(w) + \frac{\mu}{2} \|w - w^t\|_2^2
\end{equation}

To enforce sample-level differential privacy (DP), each client bounds local sample influence by clipping gradient norms to a threshold $C$. Noise calibrated via Rényi DP is added prior to local aggregation, yielding client updates $w_k^{t+1}$ that satisfy sample-level $(\epsilon, \delta)$-DP under the Gaussian mechanism. Client contribution valuation is measured using coalition utility $v(S)$, which evaluates a subset $S \subseteq \{1, \dots, K\}$ on a test set $D_{test}$.

\section{Mathematical Methodology}

\subsection{Classification Model}
The staging model is parameterized as a custom logistic regression. For a feature vector $x_i$, the log-odds are modeled as:
\begin{equation}
z_i = x_i^T w
\end{equation}
and the staging probability $p_i$ is computed using the sigmoid function:
\begin{equation}
p_i = \sigma(z_i) = \frac{1}{1 + e^{-z_i}}
\end{equation}
The local loss function $\ell(w; x_i, y_i)$ is binary cross-entropy:
\begin{equation}
\ell(w; x_i, y_i) = - [y_i \ln p_i + (1 - y_i) \ln(1 - p_i)]
\end{equation}

\subsection{Rényi Differential Privacy Accountant}
The Gaussian mechanism adds calibrated noise to local gradient updates. Let adjacent datasets $D_k, D'_k$ differ by one sample record. A randomized mechanism $\mathcal{M}$ satisfies $(\alpha, \epsilon_{RDP})$-RDP if for all adjacent $D_k, D'_k$, the Rényi divergence of order $\alpha > 1$ satisfies:
\begin{equation}
D_{\alpha}(\mathcal{M}(D_k) || \mathcal{M}(D'_k)) \le \epsilon_{RDP}
\end{equation}
At step $s$, the client clips per-sample gradients: $\bar{g}_i = g_i / \max(1, \|g_i\|_2 / C)$. The local updates are perturbed:
\begin{equation}
\tilde{g} = \frac{1}{B} \sum_{i=1}^B \bar{g}_i + z_s, \quad z_s \sim \mathcal{N}(0, \sigma^2 C^2 / B^2)
\end{equation}
For a full sampling rate $q=1.0$ (representing full-batch updates per local step), the step privacy loss is:
\begin{equation}
\epsilon_{RDP}(\alpha) = \frac{\alpha}{2 \sigma^2}
\end{equation}
After $N = T \times E$ steps, the composed RDP loss is:
\begin{equation}
\epsilon_{RDP}^{total}(\alpha) = N \cdot \frac{\alpha}{2 \sigma^2}
\end{equation}
The coordinating server converts this composed bound to standard $(\epsilon, \delta)$-DP:
\begin{equation}
\epsilon = \min_{\alpha \in \mathcal{A}} \left\{ N \cdot \frac{\alpha}{2 \sigma^2} + \frac{\ln(1/\delta)}{\alpha - 1} \right\}
\end{equation}
The server numerically solves for the noise multiplier $\sigma$ using a binary search grid over candidate Rényi orders $\mathcal{A}$.

\subsection{Client Valuation via Permutation Shapley Values}
To allocate incentive shares, we compute permutation-based Shapley values. Let $K$ be the number of hospitals, and let $S$ be a sub-coalition. The marginal value of client $k$ is computed across random permutations $\pi \in \Pi$:
\begin{equation}
\phi_k = \frac{1}{P} \sum_{j=1}^P \left[ v(S_j(\pi) \cup \{k\}) - v(S_j(\pi)) \right]
\end{equation}
where $S_j(\pi)$ is the set of clients preceding $k$ in permutation $j$, and $P$ is the number of Monte Carlo permutation samples.

\subsection{Membership Inference Attack Model}
We evaluate empirical privacy protection using a confidence-threshold Membership Inference Attack (MIA). The attacker score for sample $i$ is:
\begin{equation}
score(x_i, y_i) = y_i \cdot p_i + (1 - y_i) \cdot (1 - p_i)
\end{equation}
The attacker evaluates the ROC AUC of this score separating training members from held-out non-members.

\subsection{Domain Shift Formulations}
Robustness is evaluated under controlled synthetic shifts on the test set:
\begin{itemize}
    \item **Covariate Shift**: Scale features by factor $\beta$ and add noise: $\tilde{x} = \beta x + e, \, e \sim \mathcal{N}(0, 0.1)$.
    \item **Concept Shift**: Flip labels with probability $\gamma$: $\tilde{y} = 1-y$ with probability $\gamma$.
\end{itemize}

\section{Datasets and Preprocessing}

\subsection{TCGA-PRAD Clinical-Protein Cohort}
We utilize the primary prostate cancer cohort from TCGA-PRAD \cite{TCGA2015PRAD}. The dataset contains 572 sample records matching 500 unique patient case IDs. Because multiple biological samples originate from the same patient barcode, the differential privacy guarantee is strictly sample-level (row-level) rather than patient-level.

\subsection{Preprocessing Pipeline}
To prevent look-ahead data leakage, we partition the dataset into an 80\% training split (277 samples) and a 20\% test split (70 samples) prior to fitting any preprocessing parameters. Imputation, scaling, and categorical encoding are fitted strictly on the training partition and applied to transform the test set. Principal Component Analysis (PCA) is applied to the 456 protein expression features, retaining components explaining 95\% of the variance (115 components).

\subsection{Dirichlet Client Heterogeneity Partitioning}
To simulate non-IID clinical settings, the training split is distributed among 3 hospitals using a Dirichlet allocation over staging labels. We control the level of statistical skew by varying the concentration parameter $\alpha \in \{10.0, 1.0, 0.5, 0.1\}$, where smaller $\alpha$ indicates extreme label heterogeneity.

\subsection{Survival-Analysis Subset Limitation}
After preprocessing and cohort matching, the survival-analysis subset contains 347 observations, with 12 observed death events (9 in the training split of size 277, and 0 in the test split of size 70).
Due to the severe right-censoring in the PRAD cohort, C-index discrimination was not estimable in the test split (0 death events). No performance claims are made for survival discrimination.

\section{Experimental Protocol}
The baseline optimization parameters are set to: Rounds $T=15$, local epochs $E=3$, learning rate $\eta=0.1$, proximal parameter $\mu=0.5$, clipping norm $C=1.0$, and target privacy delta $\delta=10^{-5}$. The reproducibility details are summarized in Table \ref{tab:protocol}.

\begin{table}[ht]
\caption{Reproducible Experimental Protocol Configuration}
\centering
\begin{tabular}{ll}
\toprule
Parameter & Value \\
\midrule
Total Samples & 572 \\
Unique Patients & 500 \\
Train/Test Split & 80\% / 20\% \\
Clinical Features & 62 categorical, 12 numerical \\
PCA Proteins & 456 columns $\to$ 115 components \\
communication Rounds ($T$) & 15 \\
Local Epochs ($E$) & 3 \\
FedProx Regularizer ($\mu$) & 0.5 \\
Gradient clipping Norm ($C$) & 1.0 \\
Target Privacy delta ($\delta$) & $10^{-5}$ \\
\bottomrule
\end{tabular}
\label{tab:protocol}
\end{table}

\section{Results}

\subsection{Centralized Baselines and Staging Performance}
Centralized classification models trained on the combined training partition yield:
Lasso AUC = **0.7850**, Ridge AUC = **0.7642**, MLP NN AUC = **0.6146**, Custom NumPy Model AUC = **0.7566** [95\% Bootstrap CI: $0.6447 - 0.8659$].

\subsection{Experiment A: Bounded Client Update Drift}
We evaluated client-update divergence over 5 seeds across Dirichlet partitions $\alpha \in \{10.0, 1.0, 0.5, 0.1\}$. Enabling proximal regularization ($\mu=0.5$) consistently restricts the average client weight drift by approximately 4.2--4.3\%:
*   $\alpha = 10.0$ (Almost IID): FedAvg Drift = $0.1628 \pm 0.0119$ vs. FedProx Drift = $0.1557 \pm 0.0112$
*   $\alpha = 0.1$ (Extreme skew): FedAvg Drift = $0.1908 \pm 0.0289$ vs. FedProx Drift = $0.1825 \pm 0.0272$
This consistent containment across all heterogeneity levels supports **H1**.

\subsection{Experiment C: 2D Privacy × Heterogeneity Interaction Matrix}
We swept Dirichlet label skew ($\alpha$) and privacy budget ($\epsilon$) across 5 random seeds to evaluate global model test AUC under joint constraints in Table \ref{tab:interaction}:

\begin{table}[ht]
\caption{2D Privacy × Heterogeneity Interaction Matrix (Global Test AUC, 5 Seeds)}
\centering
\resizebox{\columnwidth}{!}{
\begin{tabular}{ccccccc}
\toprule
Dirichlet $\alpha$ & Model & $\epsilon=0.5$ & $\epsilon=1.0$ & $\epsilon=2.0$ & $\epsilon=5.0$ & $\epsilon=10.0$ \\
\midrule
10.0 (IID) & FedAvg & $0.610 \pm 0.074$ & $0.689 \pm 0.048$ & $0.741 \pm 0.017$ & $0.757 \pm 0.005$ & $0.756 \pm 0.004$ \\
           & FedProx & $0.610 \pm 0.077$ & $0.694 \pm 0.048$ & $0.742 \pm 0.017$ & $0.756 \pm 0.004$ & $0.755 \pm 0.006$ \\
           & Benefit ($\Delta$) & $-0.000 \pm 0.007$ & \textbf{$+0.004 \pm 0.002$} & \textbf{$+0.002 \pm 0.001$} & $-0.001 \pm 0.002$ & $-0.001 \pm 0.002$ \\
\midrule
1.0 (Mod.) & FedAvg & $0.611 \pm 0.075$ & $0.689 \pm 0.049$ & $0.741 \pm 0.017$ & $0.758 \pm 0.004$ & $0.756 \pm 0.005$ \\
           & FedProx & $0.611 \pm 0.077$ & $0.694 \pm 0.049$ & $0.744 \pm 0.016$ & $0.757 \pm 0.004$ & $0.755 \pm 0.006$ \\
           & Benefit ($\Delta$) & $+0.001 \pm 0.002$ & \textbf{$+0.004 \pm 0.002$} & \textbf{$+0.003 \pm 0.002$} & $-0.000 \pm 0.003$ & $-0.001 \pm 0.001$ \\
\midrule
0.5 (Str.) & FedAvg & $0.610 \pm 0.072$ & $0.691 \pm 0.047$ & $0.741 \pm 0.017$ & $0.758 \pm 0.004$ & $0.755 \pm 0.004$ \\
           & FedProx & $0.612 \pm 0.076$ & $0.694 \pm 0.047$ & $0.742 \pm 0.016$ & $0.757 \pm 0.004$ & $0.755 \pm 0.006$ \\
           & Benefit ($\Delta$) & $+0.002 \pm 0.006$ & \textbf{$+0.003 \pm 0.002$} & \textbf{$+0.002 \pm 0.001$} & $-0.001 \pm 0.002$ & $-0.000 \pm 0.001$ \\
\midrule
0.1 (Ext.) & FedAvg & $0.612 \pm 0.074$ & $0.690 \pm 0.049$ & $0.741 \pm 0.016$ & $0.758 \pm 0.004$ & $0.755 \pm 0.005$ \\
           & FedProx & $0.612 \pm 0.078$ & $0.694 \pm 0.050$ & $0.743 \pm 0.015$ & $0.757 \pm 0.004$ & $0.756 \pm 0.004$ \\
           & Benefit ($\Delta$) & $+0.001 \pm 0.005$ & \textbf{$+0.005 \pm 0.001$} & \textbf{$+0.002 \pm 0.001$} & $-0.001 \pm 0.002$ & $+0.000 \pm 0.001$ \\
\bottomrule
\end{tabular}
}
\label{tab:interaction}
\end{table}

Under joint privacy–heterogeneity constraints over repeated trials, FedProx showed a positive utility difference relative to FedAvg across the evaluated privacy–heterogeneity configurations, with the 95\% confidence intervals of the performance benefits remaining strictly positive for moderate privacy levels ($\epsilon \in \{1.0, 2.0\}$). This supports **H3**.

\subsection{Experiment D: Personalization Gains}
Evaluated client test partitions matching local Dirichlet skews ($\alpha=0.5$):
For Hospital 2 (containing valid classes), Local-Only AUC = **0.6315**, Global FedProx AUC = **0.7433**, and Personalized FL (PFL) AUC = **0.7101**. PFL fine-tuning improves staging performance by **+7.87%** ($0.7101$ vs. $0.6315$) compared to the local-only baseline, but does not outperform the global consensus model ($0.7433$) in this sparse data environment, supporting **H4**. Hospital 1 and 3 local test partitions contained only a single target class, resulting in undefined local AUCs.

\subsection{Experiment F: Shapley Valuation under Privacy Noise}
We evaluated client permutation Shapley values under varying differential privacy constraints in Table \ref{tab:shapley}:

\begin{table}[ht]
\caption{Shapley Client Scores and Ranks under Privacy Noise}
\centering
\begin{tabular}{ccccccc}
\toprule
Privacy ($\epsilon$) & Hospital 1 Score & Hospital 2 Score & Hospital 3 Score & Hospital 1 Rank & Hospital 2 Rank & Hospital 3 Rank \\
\midrule
No DP & 0.0611 & 0.0720 & 0.1065 & 3 & 2 & 1 \\
10.0  & 0.1100 & 0.0210 & 0.0944 & 1 & 3 & 2 \\
5.0   & 0.1081 & 0.0342 & 0.0906 & 1 & 3 & 2 \\
2.0   & 0.1057 & 0.0475 & 0.0807 & 1 & 3 & 2 \\
1.0   & 0.1045 & 0.0581 & 0.0713 & 1 & 3 & 2 \\
0.5   & 0.0931 & 0.0581 & 0.0505 & 1 & 2 & 3 \\
\bottomrule
\end{tabular}
\label{tab:shapley}
\end{table}

The evaluations in Table \ref{tab:shapley} demonstrate that differential privacy noise altered the observed contribution rankings in the evaluated three-client setting. Shapley values are aligned with empirical AUC degradation when removed ($\rho=0.5000$, Experiment G). For the three-client setting, the measured Shapley computation time ($0.0561s$) was comparable to LOO ($0.0438s$), although this computational advantage may not persist as the number of participating clients increases. These results support **H5**.

\subsection{Experiment I: Empirical Privacy Attack (MIA Resistance)}
Without DP, the MIA Attacker achieves an AUC of **0.5679** and an attacker advantage of **0.0734** on prediction confidence. Once DP is enabled (even at $\epsilon=10.0$), the attacker AUC drops to **0.4802** and attacker advantage falls to **-0.0161**, indicating substantially reduced vulnerability to the evaluated confidence-threshold membership-inference attack.

\subsection{Experiment H: Controlled Domain-Shift Evaluation}
We evaluated model robustness under controlled synthetic domain shifts. The test split AUC degrades gracefully from **0.7566** to **0.7462** (at severity 2.0) under covariate shift $P(X)$, but drops precipitously to **0.5265** (at severity 2.0) under concept shift $P(Y|X)$ as target labels are flipped. This distinct sensitivity supports **H6**.

\section{Survival Modeling Extension and Dataset Limitation}
After preprocessing and cohort matching, the survival-analysis subset contains **347 observations** and **12 death events** (9 in the training split of size 277, and 0 in the test split of size 70).
Although a federated Cox proportional-hazards implementation was developed, reliable out-of-sample survival discrimination could not be evaluated on the selected TCGA-PRAD split because the held-out test partition contained zero observed death events due to extreme right-censoring. Consequently, no survival-performance claim is made. This is presented as an honest dataset-level limitation regarding survival time modeling.

\section{Discussion}
The findings in this study answer our primary research question: privacy constraints and statistical heterogeneity jointly influence both optimization stability and the economics of clinical data sharing. Under extreme non-IID conditions, client parameters drift away from consensus. FedProx controls this drift by penalizing local optimization steps. Under differential privacy constraints, client updates are perturbed by noise. We observed that the FedProx regularizer provides a protective utility benefit under joint constraints, stabilizing learning paths against privacy-induced gradient noise.

Importantly, privacy mechanisms have consequences beyond model staging performance. We showed that adding differential privacy noise disrupts contribution valuation, altering client contribution rankings in the evaluated three-client setting. This suggests a conflict between privacy and incentive alignment in collaborative clinical consortia.

\section{Limitations and Future Work}
The study has several key limitations: (1) Preprocessing and DP guarantees are sample-level, not patient-level. (2) The valuation mechanism was evaluated in a three-client setting, and Shapley computational cost scales exponentially with client size. (3) Out-of-sample survival discrimination could not be evaluated due to extreme right-censoring in the PRAD cohort. Future work will investigate patient-level DP, scalable Shapley value approximations, and validation on prospective multi-institutional cohorts with sufficient survival events.

\section{Conclusion}
In this study, we systematically characterized the joint interactions between optimization drift, differential privacy boundaries, personalized adaptation, and client contribution valuation. We demonstrated that FedProx constrains client weight drift, stabilizes model performance against DP noise, and that DP noise systematically alters perceived economic contribution. Finally, we audited membership vulnerability and documented survival data censoring limitations. These multi-dimensional findings provide a mathematically rigorous blueprint for collaborative staging networks.

\bibliographystyle{IEEEtran}
\bibliography{references}

\end{document}
